% Generated by experiments/scripts/export_latex.py (export_latex_c12.py) -- do not edit by hand.
% Source: experiments/data/2-processed-tables-and-evaluations/revised-contract-12m/sybil-model/sybil_model.json
\begin{table}[htbp]
\centering
\caption[계약 그래프에서 본 시빌 비용 모형과 잔액 문턱에 대한 민감도.]{$T_{\mathrm{obs}}$의 계약 엔도스랭크 그래프에서 본 시빌 비용 모형과, 그 모형의 잔액 문턱 $\theta$에 대한 민감도예요. 그래프에는 상한을 씌운 값이 $\theta$ 이상인 허락 쌍만 남겨요. 그래서 농장 계약은 하나하나 동결 시점에 그 토큰을 $\theta$만큼 갖고 있어야 해요. 열은 차례대로 남긴 쌍과 그 몫, 그래프의 주인 수와 점 수, 매칭 코호트 점수와 주 분석 점수 사이의 순위 일치 $\tau$, 계약 10개로 된 별 모양 농장의 목표가 받는 점수를 코호트의 가장 높은 엔도스랭크에 대한 몫으로 나타낸 값, 그리고 닫힌 꼴로 계산할 때 그 가장 높은 점수에 이르는 데 필요한 농장 계약의 수예요. 농장 주소는 모두 계약이에요. 계약 하나를 만드는 데 팩토리로는 가스가 32{,}000 이상, 자기 트랜잭션으로는 53{,}000 이상 들고, 코드 한 바이트마다 가스 200이 더 들어요(EVM 수수료표의 실행 가스이고, L1에 자료를 올리는 수수료는 넣지 않았어요). 1달러가 뒷받침하는 허락의 금액 무게는 $V=1.000$이에요. 그래서 잔액이 더 커도 무게는 더 늘지 않아요. 이 값들은 어떤 규칙도 판정하지 않아요.}
\label{tab:sybil-threshold}
\fitwidth{%
\begin{tabular}{lrrrrccr}
\toprule
문턱 & 남긴 쌍 & 몫 & 주인 & 점 & 주 분석과의 $\tau$ & 별 모양 목표 / 최고 & 최고에 이르는 농장 계약 \\
\midrule
없음 (주 분석) & 114{,}997 & 100\% & 9{,}452 & 109{,}381 & 1.000 & 0.030 & 374 \\
\midrule
$\theta=\$1$ & 4{,}247 & 3.7\% & 2{,}390 & 3{,}394 & 0.709 & 0.092 & 119 \\
$\theta=\$10$ & 2{,}167 & 1.9\% & 1{,}000 & 1{,}769 & 0.589 & 0.325 & 33 \\
$\theta=\$100$ & 992 & 0.9\% & 500 & 837 & 0.464 & 1.054 & 9 \\
\bottomrule
\end{tabular}
}
\end{table}
